\documentclass[11pt]{article}

\usepackage[letterpaper,margin=1in]{geometry}
\usepackage[T1]{fontenc}
\usepackage[utf8]{inputenc}
\usepackage{lmodern}
\usepackage{microtype}
\usepackage{amsmath,amssymb,mathtools,bm}
\usepackage{enumitem}
\usepackage{xcolor}
\usepackage{booktabs}
\usepackage{array}
\usepackage{hyperref}
\usepackage[most]{tcolorbox}
\usepackage{fancyhdr}
\usepackage{titlesec}
\usepackage{parskip}

\definecolor{qftblue}{HTML}{245A8D}
\definecolor{qftgreen}{HTML}{2F6B4F}
\definecolor{qftred}{HTML}{A23B3B}
\definecolor{qftgray}{HTML}{F4F6F8}
\definecolor{qftgold}{HTML}{C08B2C}

\hypersetup{
  colorlinks=true,
  linkcolor=qftblue,
  urlcolor=qftblue,
  citecolor=qftgreen,
  pdftitle={PHYS 662 Lecture 1: From Classical Particles to Classical Fields},
  pdfauthor={Nima Lashkari}
}

\pagestyle{fancy}
\fancyhf{}
\lhead{PHYS 662: Quantum Field Theory I}
\rhead{Lecture 1}
\cfoot{\thepage}
\setlength{\headheight}{14pt}

\titleformat{\section}{\Large\bfseries\color{qftblue}}{\thesection.}{0.6em}{}
\titleformat{\subsection}{\large\bfseries\color{qftgreen}}{\thesubsection}{0.6em}{}
\titleformat{\subsubsection}{\normalsize\bfseries\color{qftred}}{\thesubsubsection}{0.6em}{}

\setlist[itemize]{leftmargin=2em,itemsep=0.25em,topsep=0.4em}
\setlist[enumerate]{leftmargin=2em,itemsep=0.25em,topsep=0.4em}

\newcommand{\dd}{\mathrm{d}}
\newcommand{\e}{\mathrm{e}}
\newcommand{\avg}[1]{\left\langle #1\right\rangle_\beta}
\newcommand{\avgh}[1]{\left\langle #1\right\rangle_{\beta,h}}
\newcommand{\avgc}[1]{\left\langle #1\right\rangle_{\beta,c}}
\newcommand{\cD}{\mathcal{D}}
\newcommand{\R}{\mathbb{R}}
\newcommand{\cO}{\mathcal{O}}
\newcommand{\cF}{\mathcal{F}}

\newtcolorbox{keyidea}[1][]{
  enhanced,
  breakable,
  colback=qftgray,
  colframe=qftblue,
  boxrule=0.8pt,
  arc=1.5mm,
  title={#1},
  fonttitle=\bfseries,
  coltitle=qftblue
}

\newtcolorbox{remarkbox}[1][]{
  enhanced,
  breakable,
  colback=white,
  colframe=qftgreen,
  boxrule=0.7pt,
  arc=1.5mm,
  title={#1},
  fonttitle=\bfseries,
  coltitle=qftgreen
}

\newtcolorbox{warningbox}[1][]{
  enhanced,
  breakable,
  colback=white,
  colframe=qftred,
  boxrule=0.7pt,
  arc=1.5mm,
  title={#1},
  fonttitle=\bfseries,
  coltitle=qftred
}

\begin{document}
\pagenumbering{gobble}

\begin{titlepage}
\centering
\vspace*{1.2in}
{\scshape\Large PHYS 662 --- Quantum Field Theory I\par}
\vspace{0.35in}
{\Huge\bfseries\color{qftblue} Lecture 1\par}
\vspace{0.15in}
{\LARGE\bfseries From Classical Particles to Classical Fields\par}
\vspace{0.12in}
{\Large\color{qftgreen} Statistical Field Theory as a Long-Distance Description\par}
\vspace{0.6in}
{\large Nima Lashkari\par}
{\large Fall 2026\par}
\vfill
\begin{minipage}{0.86\textwidth}
\small
Statistical field theory can emerge for the effective low-energy dynamics of a classical many-body problem with local interactions, when one asks only about physics at distances much larger than the microscopic cutoff. The central ideas that underlie the emergence of (statistical) field theory are locality, symmetry, (thermal) fluctuations and coarse-graining.

\medskip
Course page: \url{https://nimalashkari.com/quantum-field-theory-2026-part-i/}
\end{minipage}
\vspace{0.5in}
\end{titlepage}

\pagenumbering{roman}
\tableofcontents
\newpage
\pagenumbering{arabic}

\section*{Course perspective}
\addcontentsline{toc}{section}{Course perspective}

In this course, we will assume fluency in quantum mechanics and statistical mechanics, and familiarity with special relativity and basic group theory.  There is no single textbook that we will follow.  The course website is the authoritative source for the syllabus, homework, lecture notes, and recordings.

The first part of the course develops \emph{classical field theory}; the second part quantizes it, first canonically and then using path integrals.  Along the way, three principles will repeatedly organize the subject:
\begin{center}
\begin{tabular}{>{\bfseries}p{0.20\textwidth} p{0.67\textwidth}}
Fluctuations & Correlation functions, probability distributions, and eventually path integrals. \\
Locality & Local interactions, local fields, causality, and the action principle. \\
Symmetry & Constraints on dynamics, conservation laws, charges, and spin. \\
\end{tabular}
\end{center}
A fourth principle, which will become increasingly important, is \emph{renormalization}: the long-distance theory often forgets most microscopic details.

\begin{keyidea}[A useful attitude]
Quantum field theory is not a finished collection of recipes.  It is a framework whose conceptual structure is still being developed.  Questions are therefore not interruptions to the subject; they are part of the subject.  The goal of this course is to learn the standard tools while keeping track of what they mean and why they work.
\end{keyidea}

\section{Why start with statistical field theory?}

Field theory is a language for systems whose relevant degrees of freedom are functions of position (and, in dynamical problems, time).  It is useful both at very high energies and at very low energies.  In high-energy physics, relativistic quantum field theory is our framework for elementary interactions.  In many-body physics, fields emerge as collective variables describing long-wavelength fluctuations of a large number of microscopic degrees of freedom.

Particles are extremely useful when a theory has stable or sufficiently long-lived, well-separated excitations.  But a particle description is not always available.  Near a continuous phase transition, for example, fluctuations occur on arbitrarily large length scales, and the natural observables are correlation functions of coarse-grained fields.  Similarly, in relativistic QFT, particle states are best thought of as a special asymptotic or low-energy organization of the theory rather than as the only possible starting point.

\begin{remarkbox}[Fields versus particles]
It is often said that ``fields are more fundamental than particles.''  The precise statement we will use is more modest: \emph{the local field-theoretic description is often more robust than a particle description}.  A theory may have no useful particle interpretation in some regime, while its local observables and correlation functions remain perfectly meaningful.
\end{remarkbox}

Classical statistical field theory is the cleanest place to see how such a description emerges.  It will also give us our first example of a path integral.

\section{Thermal fluctuations in a finite system}

Consider first a finite classical system with phase-space coordinates
\begin{equation}
  \Gamma=(q_1,\ldots,q_N,p_1,\ldots,p_N)
\end{equation}
and Hamiltonian $H(\Gamma)$.  A microstate is a point in phase space.  If the system is weakly coupled to a thermal bath at inverse temperature
\begin{equation}
  \beta=\frac{1}{k_B T},
\end{equation}
then we no longer assign a single microstate to the system.  Instead we assign the canonical probability measure
\begin{equation}
  \dd\mu_\beta(\Gamma)=\frac{1}{Z(\beta)}\,\e^{-\beta H(\Gamma)}\,\dd\Gamma,
  \qquad
  Z(\beta)=\int \dd\Gamma\,\e^{-\beta H(\Gamma)}.
  \label{eq:canonical}
\end{equation}
For discrete degrees of freedom the phase-space integral is replaced by a sum.  We will suppress conventional normalization factors in $\dd\Gamma$ that do not affect the discussion.

For an observable $A(\Gamma)$,
\begin{equation}
  \avg{A}=\frac{1}{Z}\int \dd\Gamma\,\e^{-\beta H(\Gamma)}A(\Gamma).
\end{equation}
The probability distribution turns deterministic microscopic variables into fluctuating statistical observables.

\subsection{Correlation functions and connected correlators}

For two observables $A$ and $B$, the simplest measure of correlated fluctuations is
\begin{equation}
  \avgc{AB}
  \equiv
  \avg{AB}-\avg{A}\avg{B}.
  \label{eq:connected}
\end{equation}
For example,
\begin{equation}
  \avg{(q_1-\avg{q_1})(q_2-\avg{q_2})}
  =\avg{q_1q_2}-\avg{q_1}\avg{q_2}.
\end{equation}
If the joint distribution factorizes, $p(q_1,q_2)=p(q_1)p(q_2)$, then this connected two-point function vanishes.  Nonzero connected correlators therefore measure the failure of statistical independence.

Higher moments and connected moments are the basic observables of statistical field theory and, later, of quantum field theory.

\subsection{The partition function as a generating function}

Introduce sources $J_i$ and $K_i$ coupled to positions and momenta:
\begin{equation}
  Z[J,K]
  =\int \dd\Gamma\,
  \exp\left[-\beta H(\Gamma)+\sum_i J_i q_i+\sum_i K_i p_i\right].
  \label{eq:sourceZ}
\end{equation}
Differentiating $Z$ generates ordinary moments.  For instance,
\begin{equation}
  \left.\frac{1}{Z}\frac{\partial Z}{\partial J_i}\right|_{J=K=0}=\avg{q_i},
  \qquad
  \left.\frac{1}{Z}\frac{\partial^2 Z}{\partial J_i\partial J_j}\right|_{0}=\avg{q_iq_j}.
\end{equation}
Even more importantly,
\begin{equation}
  W[J,K]\equiv \log Z[J,K]
\end{equation}
generates \emph{connected} correlation functions:
\begin{equation}
  \left.\frac{\partial^2 W}{\partial J_i\partial J_j}\right|_{0}
  =\avgc{q_iq_j}.
\end{equation}
This distinction between $Z$ and $W=\log Z$ will reappear throughout QFT.

\subsection{Fluctuations as response}

There is a direct physical meaning to connected correlators.  Perturb the Hamiltonian by a source $h$ coupled to an observable $A$,
\begin{equation}
  H_h=H-hA.
\end{equation}
Then the response of another observable $B$ is
\begin{equation}
  \left.\frac{\partial}{\partial h}\avgh{B}\right|_{h=0}
  =\beta\avgc{AB}.
  \label{eq:staticFDT}
\end{equation}
For $A=B$, the static susceptibility is proportional to the variance.  Equation~\eqref{eq:staticFDT} is the simplest form of the fluctuation--dissipation relation: fluctuations in equilibrium determine linear response to a weak external perturbation.

\section{A laboratory: the \texorpdfstring{$O(N)$}{O(N)} spin model}

A useful family of many-body systems consists of spins or rotors living on the sites of a lattice.  At each site $i$, let
\begin{equation}
  \bm S_i\in \R^N,
  \qquad
  \bm S_i\cdot \bm S_i=1.
\end{equation}
The nearest-neighbor ferromagnetic Hamiltonian is
\begin{equation}
  H=-J\sum_{\langle ij\rangle}\bm S_i\cdot\bm S_j,
  \qquad J>0.
  \label{eq:ONH}
\end{equation}
The model is invariant under a common orthogonal rotation of all spins,
\begin{equation}
  \bm S_i\longrightarrow R\bm S_i,
  \qquad
  R^T R=\mathbf 1,
\end{equation}
so its global symmetry group is $O(N)$.

Important special cases are
\begin{center}
\begin{tabular}{ccl}
\toprule
$N$ & Degree of freedom & Model \\
\midrule
$1$ & $S_i=\pm1$ & Ising model \\
$2$ & $\bm S_i=(\cos\theta_i,\sin\theta_i)$ & XY model \\
$3$ & $\bm S_i\in S^2$ & Heisenberg model \\
\bottomrule
\end{tabular}
\end{center}
For $N=1$, $O(1)\simeq \mathbb Z_2$ and the degree of freedom is discrete.  For $N\ge2$ the spin takes values on the sphere $S^{N-1}$.

The thermal state is
\begin{equation}
  p(\{\bm S_i\})=\frac{1}{Z}\e^{-\beta H(\{\bm S_i\})}.
\end{equation}
At nonzero temperature, the microscopic spins fluctuate.  What survives at long distances are suitable averages and correlation functions.

\subsection{Quasiparticles are emergent}

Even a spin system may admit a particle-like description at sufficiently low energy.  A simple example is the one-dimensional Ising chain,
\begin{equation}
  H=-J\sum_i s_i s_{i+1},
  \qquad s_i=\pm1.
\end{equation}
The ground states have all spins aligned.  A bond with $s_i\neq s_{i+1}$ is a \emph{domain wall}.  Relative to an aligned bond, one domain wall costs energy $2J$.  With periodic boundary conditions domain walls are created in pairs.  Under an appropriate spin-flip dynamics they can move, and two domain walls can annihilate when they meet.

Thus the number of these particle-like excitations is not a microscopic conserved quantity.  The ``particles'' are defects of a more fundamental many-body configuration.  This is our first example of the general QFT idea that particle number need not be fundamental.

\section{Correlation length}

For a translationally invariant spin system, consider the connected two-point function
\begin{equation}
  C(\bm r)
  =\avg{\bm S_{\bm 0}\cdot\bm S_{\bm r}}
   -\avg{\bm S_{\bm 0}}\cdot\avg{\bm S_{\bm r}}.
  \label{eq:spinC}
\end{equation}
In a homogeneous state this depends only on the separation $r=|\bm r|$ and on thermodynamic parameters such as temperature.

If the Gibbs state is in a clustering phase, then
\begin{equation}
  C(r)\longrightarrow 0,
  \qquad r\longrightarrow\infty.
  \label{eq:cluster}
\end{equation}
The important question is \emph{how quickly} it approaches zero.

Away from criticality, one generically expects an exponential length scale,
\begin{equation}
  C(r)\sim \frac{1}{r^{p}}\e^{-r/\xi},
  \qquad r\gg a,
  \label{eq:massivecorr}
\end{equation}
where $a$ is the microscopic lattice spacing and $\xi$ is the \emph{correlation length}.  The power $p$ depends on the theory and is less important here than the exponential cutoff.  In field-theory language, $\xi^{-1}$ plays the role of a mass scale.

At a scale-invariant critical point the correlation length diverges, and the exponential disappears:
\begin{equation}
  C(r)\sim \frac{1}{r^{2\Delta_S}}.
  \label{eq:criticalcorr}
\end{equation}
Near a continuous phase transition,
\begin{equation}
  \xi\sim |t|^{-\nu},
  \qquad
  t\equiv \frac{T-T_c}{T_c},
  \label{eq:xinu}
\end{equation}
with critical exponent $\nu>0$.

\begin{warningbox}[A clustering subtlety]
Equation~\eqref{eq:cluster} is naturally stated for a pure thermodynamic phase (more precisely, an extremal Gibbs state).  A symmetric mixture of two symmetry-broken phases need not cluster even though each pure phase does.  Connected correlators remove the one-point function, but phase mixtures can still retain long-range correlations.
\end{warningbox}

\subsection{Physical meaning of \texorpdfstring{$\xi$}{xi}}

Two variables separated by a distance much larger than $\xi$ are approximately statistically independent for many long-distance questions.  This suggests dividing the system into regions of linear size of order $\xi$: the microscopic variables inside a region are strongly correlated, while widely separated regions are only weakly correlated.

This picture works well when $\xi$ is finite.  At a continuous phase transition, however, $\xi$ diverges.  There is then no finite block size beyond which all fluctuations become independent.  Fluctuations survive on every scale.  This is precisely the regime where a continuum field theory becomes natural.

\section{Separation of scales, coarse-graining, and universality}

Let $a$ denote the microscopic lattice spacing.  The meaningful continuum parameter is not simply $\xi$, but the dimensionless ratio
\begin{equation}
  \frac{\xi}{a}.
\end{equation}
A continuum field theory is appropriate when
\begin{equation}
  \boxed{\frac{\xi}{a}\gg 1.}
  \label{eq:continuumcriterion}
\end{equation}
Then there is a large window of scales
\begin{equation}
  a\ll \ell\ll \xi
\end{equation}
on which the microscopic lattice is invisible while correlations are still nontrivial.

We can coarse-grain the system by grouping microscopic spins into blocks of size $\ell$.  If $B_X$ is a block centered near the coarse-grained point $X$, define the block magnetization
\begin{equation}
  m_X=\frac{1}{N_B}\sum_{i\in B_X}s_i.
  \label{eq:blockspin}
\end{equation}
Although each $s_i$ takes only the values $\pm1$, the average $m_X$ takes many closely spaced values when $N_B\gg1$.  It can therefore be approximated by a continuous variable.

At distances much larger than $a$, many different microscopic Hamiltonians can flow to the same long-distance description.  Their detailed short-distance correlations differ, but their large-scale observables agree after a finite number of parameters have been matched.  This is \emph{universality}.

\begin{keyidea}[Wilsonian viewpoint]
Coarse-graining generates every interaction compatible with the symmetries, but most of these interactions become less important at long distances.  Near a renormalization-group fixed point, only a small set of relevant or marginal couplings control the infrared physics.  In this sense, long-distance physics ``forgets'' most microscopic details.
\end{keyidea}

The word \emph{renormalizability} will later acquire a more precise meaning.  For now, the important principle is the separation of scales: short-distance structure can influence a small number of effective couplings without forcing us to keep track of every microscopic variable.

\section{From spins to a field}

We now make the passage from a many-body system to a statistical field theory explicit.

\subsection{Reorganizing the partition function}

Start from the Ising partition function
\begin{equation}
  Z=\sum_{\{s_i=\pm1\}}\e^{-\beta H[\{s_i\}]}.
\end{equation}
Suppose we group microscopic spin configurations according to a macroscopic variable $m$.  For a single global magnetization,
\begin{equation}
  m=\frac{1}{N}\sum_{i=1}^N s_i,
\end{equation}
we may write the exact identity
\begin{equation}
  Z
  =\sum_m\sum_{\{s_i\}:\,m(\{s_i\})=m}\e^{-\beta H[\{s_i\}]}
  \equiv
  \sum_m \e^{-\beta F(m)}.
  \label{eq:globalF}
\end{equation}
This defines a constrained free energy $F(m)$.  The free energy contains both energetic and entropic information: it counts, with Boltzmann weights, all microscopic configurations compatible with the same macroscopic magnetization.

The more useful construction is local.  Divide the lattice into blocks and keep the block magnetization $m_X$ of each block.  Then
\begin{equation}
  \e^{-\beta F[\{m_X\}]}
  \equiv
  \sum_{\{s_i\}\,\text{consistent with }\{m_X\}}
  \e^{-\beta H[\{s_i\}]},
  \label{eq:constrainedF}
\end{equation}
so that
\begin{equation}
  Z=\sum_{\{m_X\}}\e^{-\beta F[\{m_X\}]}.
  \label{eq:latticefunctional}
\end{equation}
The microscopic variables have been integrated out, leaving an effective theory of coarse-grained variables.

\subsection{The functional integral}

When the blocks contain many spins, $m_X$ is approximately continuous.  When the coarse lattice spacing is also small compared with the physical scales of interest, the collection $\{m_X\}$ becomes a smooth field $m(\bm x)$.  The discrete sum in \eqref{eq:latticefunctional} becomes formally
\begin{equation}
  \boxed{
  Z=\int \cD m\,\e^{-\beta F[m]}.
  }
  \label{eq:functionalintegral}
\end{equation}
This is a classical statistical \emph{path integral} or \emph{functional integral}.  The probability distribution is now a probability distribution on functions,
\begin{equation}
  \mathcal P[m]=\frac{1}{Z}\e^{-\beta F[m]}.
\end{equation}
The symbol $\cD m$ is not mysterious: on the lattice it is simply an ordinary product of sums or integrals over the variables $m_X$.  The continuum notation is shorthand for the limit in which the lattice regulator is taken much smaller than all distances we wish to resolve.

\section{Why the effective free energy is local}

For a short-range microscopic Hamiltonian and a slowly varying coarse-grained field, the effective free energy can be organized in a derivative expansion.  The simplest way to see the origin of a gradient term is to rewrite the Ising nearest-neighbor interaction.  Since $s_x^2=1$,
\begin{equation}
  -J s_x s_{x+a}
  =\frac{J}{2}(s_{x+a}-s_x)^2-J.
  \label{eq:isinggradient}
\end{equation}
The additive constant does not affect normalized probabilities.  Replacing the discrete variable by a smooth field gives
\begin{equation}
  \phi(x+a)-\phi(x)=a\,\partial_x\phi(x)+\cO(a^2),
\end{equation}
so the lattice difference produces a continuum gradient term $(\partial\phi)^2$.

To mimic an Ising variable by a continuous ``soft spin,'' one may introduce a local double-well potential
\begin{equation}
  V(\phi)=g(\phi^2-1)^2,
  \qquad g\gg1,
\end{equation}
which strongly favors $\phi\approx\pm1$.  After coarse-graining, however, the most general local free-energy functional should be written in terms of all operators allowed by symmetry.  For a $\mathbb Z_2$-symmetric scalar order parameter,
\begin{equation}
  \beta F[\phi]
  =\int \dd^d x\,
  \left[
  \frac{\kappa}{2}(\nabla\phi)^2
  +\frac{r}{2}\phi^2
  +\frac{u}{4!}\phi^4
  +c_6\phi^6
  +c_{2,4}\phi^2(\nabla\phi)^2
  +\cdots
  \right].
  \label{eq:LG}
\end{equation}
The ellipsis contains higher powers and higher derivatives.  The symmetry $\phi\to-\phi$ forbids odd powers in zero external field.

For the $O(N)$ model the scalar field is replaced by an $N$-component field $\bm\phi$, and the local potential depends on $\bm\phi^2$:
\begin{equation}
  \beta F[\bm\phi]
  =\int \dd^d x\,
  \left[
  \frac{\kappa}{2}(\nabla\bm\phi)^2
  +\frac{r}{2}\bm\phi^2
  +\frac{u}{4!}(\bm\phi^2)^2
  +\cdots
  \right].
  \label{eq:ONLG}
\end{equation}
This is the Landau--Ginzburg--Wilson description of the long-distance fluctuations.

\begin{remarkbox}[Where locality enters]
Locality does not mean that the exact coarse-grained functional contains only finitely many terms.  Rather, for short-range microscopic interactions and fields varying slowly compared with the cutoff, nonlocal dependence can be expanded in local operators with increasing numbers of derivatives.  Renormalization tells us which of these infinitely many local terms matter at long distances.
\end{remarkbox}

\section{The continuum limit and criticality}

Equation~\eqref{eq:continuumcriterion} is the central criterion for a continuum description.  If $\xi/a$ remains finite, then a sufficiently coarse observer eventually sees almost independent blocks, and the microscopic cutoff remains physically relevant.  To obtain a nontrivial continuum theory while sending $a\to0$, one tunes parameters so that the correlation length in lattice units diverges:
\begin{equation}
  \frac{\xi}{a}\longrightarrow\infty.
\end{equation}
A continuous phase transition provides precisely such a limit.

At criticality there is no preferred finite correlation length.  Correlation functions become scale-free, and the system exhibits fluctuations at all scales.  The continuum field theory is therefore not a description of one special length scale; it is the description that remains after the microscopic scale has become irrelevant.

\begin{warningbox}[Dimensionality]
A classical statistical system defined on a $d$-dimensional spatial lattice gives, in the continuum, a \emph{$d$-dimensional Euclidean statistical field theory}.  The extra Euclidean time dimension appears when a \emph{$d$-dimensional quantum system} is represented by an imaginary-time path integral, producing a $(d+1)$-dimensional Euclidean field theory.  These are related ideas, but they should not be conflated.
\end{warningbox}

\section{Four structural lessons}

The discussion can be summarized by four principles that will guide the rest of the course.

\begin{enumerate}
  \item \textbf{Fluctuations.}  A thermal state is a probability distribution, and its physical content is encoded in correlation functions.  The generating functional packages all such correlations into one object.

  \item \textbf{Locality.}  Short-range microscopic interactions and slowly varying fields lead to a local derivative expansion of the effective action or free energy.

  \item \textbf{Symmetry.}  Symmetries constrain which terms may appear in the effective theory.  They organize observables and, in dynamical theories, lead to conserved charges.

  \item \textbf{Renormalization and universality.}  Coarse-graining changes the couplings but often leaves long-distance observables insensitive to most microscopic details.  A continuum field theory captures the stable information that survives repeated coarse-graining.
\end{enumerate}

These principles are already visible in classical statistical mechanics.  Quantum field theory will add quantum fluctuations, relativistic causality, operator-valued observables, and the possibility of particle creation and annihilation, but the structural ideas are closely related.

\section{Connection with the law of large numbers and the central limit theorem}

There is a useful probabilistic way to think about the emergence of field theory.  Suppose we average $N$ independent or weakly correlated microscopic variables with finite variance.  The law of large numbers says that the average becomes sharply concentrated as $N\to\infty$, while the central limit theorem says that the properly rescaled fluctuations approach a Gaussian distribution.

Away from a critical point, a system with finite correlation length can be divided into blocks much larger than $\xi$.  Those blocks are approximately independent, so a central-limit mechanism often drives sufficiently coarse observables toward Gaussian statistics.  The Gaussian field theory is, in this sense, a natural fixed point of coarse-graining.

At a critical point, however, $\xi=\infty$.  Blocks at every scale remain correlated, so the assumptions behind the naive central limit theorem fail.  Coarse-graining may then approach a non-Gaussian fixed point characterized by nontrivial critical exponents.  Wilsonian renormalization can be viewed as a far-reaching generalization of the central-limit idea: instead of asking only for the limiting distribution of sums of independent variables, we classify the possible fixed points of repeated coarse-graining in interacting many-body systems.

\begin{keyidea}[Preview]
A free field theory is the field-theoretic analogue of a Gaussian distribution.  Interacting critical field theories are the non-Gaussian fixed points that survive when correlations remain important at every scale.
\end{keyidea}

\section{Questions to think about}

\begin{enumerate}
  \item What precisely is a quasiparticle, and how is it different from a fundamental microscopic degree of freedom?
  \item Does the existence of quasiparticles imply that fields are more fundamental than particles, or only that the particle description is regime-dependent?
  \item Can a field theory have no particle interpretation at all?
  \item Starting from a field theory, what additional assumptions are required before it makes sense to define particles?
  \item In what precise sense is the renormalization group a generalization of the central limit theorem?
\end{enumerate}

\enlargethispage{2\baselineskip}
\section*{Summary}
\addcontentsline{toc}{section}{Summary}

A thermal many-body system is described by a probability distribution whose connected correlation functions reveal a correlation length $\xi$.  When $\xi/a\gg1$, block variables become approximately continuous fields and the microscopic partition sum becomes a functional integral, $Z=\int\cD\phi\,\e^{-\beta F[\phi]}$.  Locality organizes $F$ in a derivative expansion, symmetry restricts the allowed terms, and renormalization explains why distinct microscopic systems can share the same continuum physics.  This is the basic route from classical spins or particles to classical fields; the next lectures ask how the story changes for quantum systems.

\end{document}
